\chapter{Preparing data}
\label{ch:data}

\section{Directory layout}

One directory per material:

\begin{lstlisting}
data/vasp/
    struct_000/   POSCAR  INCAR  POTCAR  CHGCAR  TAUCAR
    struct_001/   POSCAR  INCAR  POTCAR  CHGCAR  TAUCAR
    ...
\end{lstlisting}

The prefix is configurable (\opt{data.pattern}); anything matching it and
containing the required files is picked up.

\begin{ptip}
\textbf{Any standard VASP output works.} Poraqu\^e computes the external ionic
potential natively from \file{POSCAR}, \file{INCAR} and \file{POTCAR} using the
tabulated pseudopotential, matching a reference potential to a relative
$5\times10^{-5}$. Only \file{CHGCAR} and \file{TAUCAR} are needed as targets.
\end{ptip}

\section{Grids may differ between materials}

They usually do: \opt{ENCUT} and the cell shape fix \opt{NGXF, NGYF, NGZF}, so
two structures of the same compound can land on different meshes. This is
handled --- batches are grouped by grid shape and one model serves them all.

Fields are reduced to a common working resolution (\opt{data.resolution}, the
longest axis) by Fourier truncation, which is the exact band-limited projection
for a plane-wave field: periodicity and the electron count survive to machine
precision.

\section{Checking your data}

The sharpest check on a \file{CHGCAR} is its integral, which must come back as
the exact electron count $\sum_a Z^{\mathrm{val}}_a$:

\begin{lstlisting}
from poraque.fields import ChargeDensity, FieldGrid

path = "data/vasp/struct_000/CHGCAR"
rho = ChargeDensity.read(path, grid=FieldGrid.from_file(path))
print(rho.integrate())   # 297.0 for 27 Pt atoms at ZVAL = 11
\end{lstlisting}

A result that misses the integer points at a truncated or misread file long
before any model is trained.

\opt{poraque-train} does not repeat this check, but it does stop with an error
on a missing \file{CHGCAR}, and it reports how many points a field drives
negative when it is band-limited to \opt{data.resolution} --- Gibbs ringing
around the core peaks, which is an artefact of the truncation rather than a
fault in the data.

A missing \file{TAUCAR} is \emph{not} an error. It simply means the material
cannot serve \opt{chg2tau}; it remains a perfectly good \opt{ext2chg} training
example, the cache logs \opt{[no TAUCAR]} beside it, and a run asking for
\opt{task: all} trains the task the data can support and reports the other as
skipped.

\section{Downloading from the Materials Project}
\label{sec:mp-download}

Installing Poraqu\^e registers \opt{poraque-mp}, which turns a \emph{chemical
space} --- a set of elements --- into a local dataset of charge densities.
Every material whose composition uses only those elements is fetched, across
all stoichiometries and structures.

Size it before committing to the transfer. Over a chemical space the estimate
is exact rather than modelled: charge densities are objects in a public S3
bucket, so their sizes are read with \opt{HEAD} requests that move no payload.

\begin{lstlisting}
poraque-mp --elements Pt Pd Ni --estimate
\end{lstlisting}

\begin{lstlisting}
==================================================================
  CHGCAR ESTIMATE - Pt-Pd-Ni space
==================================================================
  Materials in space           : 27
  FILES TO DOWNLOAD            : 27
  TOTAL DOWNLOAD               : 193.3 MB
  Storage if kept gzipped      : 169.9 MB
  ...
  27 file(s) to download, 193.3 MB total (169.9 MB on disk gzipped).
  Dry run: nothing downloaded, nothing written.
\end{lstlisting}

\opt{--estimate} is a \textbf{pure dry run}: it reports to the console and
leaves no file behind --- no \file{summary.csv}, no \file{summary.json}, no
\file{chgcar\_estimate.csv}, not even a directory. Asking how big something
would be should not create anything.

Then download. \opt{--output} (equivalently \opt{--outdir}) defaults to the
\textbf{current directory}, so a command that writes hundreds of megabytes puts
them where you ran it rather than somewhere it chose:

\begin{lstlisting}
poraque-mp --elements Pt Pd Ni --output data/MP --max-size-mb 20
\end{lstlisting}

A download is \textbf{one directory per material, named by its id} --- the
shape of a VASP run, and the shape \opt{data.data\_paths} already reads:

\begin{lstlisting}
data/MP/
    summary.csv  manifest.json  manifest.csv
    structures/mp-124.cif
    mp-124/  CHGCAR.gz  INCAR  KPOINTS  POSCAR  mp.json
    mp-81/   CHGCAR.gz  INCAR  KPOINTS  POSCAR  mp.json
    mp-126/  fields.h5  ...                       # with --hdf5
\end{lstlisting}

The three VASP inputs reconstruct the calculation, so a material directory is
one VASP can be pointed at. \file{INCAR} and \file{KPOINTS} come from the task
record the density belongs to; MP's standard static set fills in when a record
cannot be read, and the file says which of the two it is. The \file{POSCAR}
comes from the \textbf{density's own header} --- that is the geometry it was
computed at, and taking it from anywhere else is how a directory ends up with a
structure and a density that disagree. \opt{--no-vasp-inputs} skips all three.

\begin{ptip}
\textbf{None of them is training data.} The loader records the density and
nothing else, and the external potential stays \emph{computed} --- exactly as
it must be, since at inference time there is no DFT run to read one from. The
\file{mp.json} beside them is what keeps the directory from being taken for a
VASP run now that it looks like one; deleting it changes which source reads the
data.
\end{ptip}

Useful filters: \opt{--band-gap MIN MAX} (\opt{0 0} selects metals),
\opt{--crystal-system}, \opt{--stable-only}, \opt{--num-sites MIN MAX}, and
\opt{--max-sites} / \opt{--max-size-mb} / \opt{--limit} to cap the download.
The size cap has \textbf{two spellings and they are the same cap}:
\opt{--max-size-gb 0.02} is \opt{--max-size-mb 20}, a decimal factor of 1000
apart, matching every size the report prints. A single \file{CHGCAR} is an
MB-scale object and a whole download is a GB-scale one, so the unit that reads
naturally depends on which of the two you are thinking about; passing both is
refused rather than resolved.
The run is resumable and one failure never aborts the batch; \file{manifest.csv}
records what landed, what was cached and what failed.

The API key is read from \opt{\$MP\_API\_KEY}, then a local \file{.env}, then
\file{\textasciitilde/.env}, so it never has to appear in a config file or a
shell history.

\begin{ptip}
\textbf{Leave the files gzipped.} Poraqu\^e reads \file{.gz}, \file{.bz2},
\file{.xz} and \file{.zip} volumetric files in place, streaming them a line at
a time. Expanding them buys nothing and costs roughly three times the storage.
\end{ptip}

\section{The whole database, not a chemical space}
\label{sec:mp-all}

\opt{--all} asks the other question: every material the index says has a charge
density, whatever it is made of. It is a single server-side query ---
\opt{has\_props=["charge\_density"]} on the summary endpoint --- rather than one
query per chemical system, and rather than pulling every summary document in
the database and discarding the ones without a density.

\begin{lstlisting}
poraque-mp --all --estimate --sample 20   # what would 20 of them cost?
poraque-mp --all --sample 20 --output data/MP    # fetch exactly those 20
\end{lstlisting}

\begin{pwarn}
\textbf{It has to be written out.} Omitting \opt{--elements} does not mean
\emph{everything}: a command line naming neither it nor \opt{--all} is refused.
That flag was required for most of this tool's life, and reading its absence as
the whole database would turn a forgotten argument into a multi-terabyte
transfer.
\end{pwarn}

Every filter from the previous section still applies, so \opt{--all} is also
how a \emph{subset} of the whole archive is chosen --- not 150\,000 materials
but the ones small enough to be worth training on:

\begin{lstlisting}
poraque-mp --all --num-sites 1 12 --max-size-mb 20 \
    --output data/MP --hdf5 --compression gzip
\end{lstlisting}

\subsection*{\opt{--hdf5} keeps everything, and VASP still cannot read it}

\opt{--hdf5} writes each density as a chunked, compressed field store rather
than a \file{CHGCAR}. It is a re-encoding, not a conversion --- pymatgen
already holds $\rho\Omega$, which is the convention the store keeps --- and it
carries \textbf{everything the text file does}: the magnetisation block of a
spin-polarised density, and the PAW augmentation records, the one-centre terms
\opt{ICHARG = 1} will not restart without.

\begin{pwarn}
\textbf{A store written before 2026-09-01 has neither.} It kept the total block
alone and wrote no records, so one download produced two different datasets:
\opt{data.spin: auto} resolves against what is on disk, and the same materials
trained a two-channel operator as \file{CHGCAR}s and a one-channel one as
stores. Re-fetch the material if you need what the old store dropped.
\end{pwarn}

What no store can do is be opened by VASP, so a run needs the density written
back out:

\begin{lstlisting}
poraque-vasp chgcar data/MP/mp-124/fields.h5 --output CHGCAR

# or straight into a deck, which converts rather than copies
poraque-vasp energy data/MP/mp-124/fields.h5 --like <run> --copy-density
\end{lstlisting}

Every mode of \opt{poraque-vasp} takes a store where it takes a \file{CHGCAR};
address one field as \opt{fields.h5::CHGCAR} when the store holds several. A
text source is copied byte for byte rather than parsed and re-emitted ---
re-rendering a file that was already correct can only lose digits or shift a
column, and the density block is a column-positional Fortran read.

\subsection*{A sample is a selection, not just a sizing trick}

\opt{--sample N} draws \opt{N} materials at random and \textbf{those become the
working set}. With \opt{--estimate} it sizes exactly those \opt{N}; without
\opt{--estimate} it downloads them. Everything the run writes into
\opt{--output} is about that subset --- \file{summary.csv} and the CIFs
included --- so a sampled directory never carries an index of a set it does not
hold.

\begin{pwarn}
\textbf{Before 2026-09-01 \opt{--sample} reached the dry run alone.} Used
without \opt{--estimate} it was silently ignored, and the command fetched the
whole matching set --- with \opt{--all}, the entire database.
\end{pwarn}

Sizing changes with it. Over a chemical space every object is \opt{HEAD}ed and
the total is exact; over the database that is tens of thousands of requests for
a number nobody needs to four digits. A sampled estimate measures its \opt{N}
exactly --- so \opt{TOTAL DOWNLOAD} is a measurement of what a fetch would
actually transfer, not an extrapolation --- and \emph{also} projects to the
whole set, as the sample mean times the number advertised, on its own line and
labelled as a projection. Both are printed because they answer different
questions, \emph{what will this fetch} and \emph{what would all of it cost},
and a report stating one of them gets read as the other.

\opt{--seed} decides \emph{which} subset, and that is what makes a sampled
download resumable: the same seed selects the same materials, so an interrupted
run resumes onto the same subset rather than fetching \opt{N} more beside it.
Sizes are also strongly right-tailed, so two draws of twenty can differ by tens
of percent, and a projection that moved every time it was asked could not be
planned against.

\section{What an MP download does and does not contain}

It contains one compressed \file{CHGCAR} per material directory and nothing
else --- no \file{POSCAR}, \file{INCAR}, \file{POTCAR}, \file{OUTCAR} or
\file{TAUCAR}. (That emptiness is also what tells the directory apart from a
run whose inputs are there to build $\vext$ from exactly.)
Three consequences, and Poraqu\^e handles the first two for you:

\begin{description}
  \item[The structure is recoverable.] A \file{CHGCAR} carries its own
    \file{POSCAR} in its first lines, so the geometry comes from the density
    itself.
  \item[The valence charges are recoverable.] A pseudopotential \file{CHGCAR}
    integrates to its cell's valence electron count, which is one linear
    equation per material in the per-element $Z^{\mathrm{val}}$. A chemical
    space of any breadth over-determines them. On the Ag--Pt--Pt set this
    returns $11$, $11$ and $10$ --- the \file{POTCAR} values, read off the
    data, from three small files.
  \item[$\tau$ is not.] No public archive publishes a kinetic energy density,
    so \opt{chg2tau} cannot be trained on MP data at all. Use
    \opt{task: ext2chg}.
\end{description}

\section{The external potential without a POTCAR}
\label{sec:mp-potcar-dir}

The one thing the density cannot supply is the \emph{pseudopotentials}, and
those are what the exact external potential is built from. Since $\vext$ is the
model's \emph{input}, this is the single most consequential setting for a run
on archive data.

Point \opt{data.potcar\_dir} at a \file{POTCAR} library --- one subdirectory
per species, as VASP distributes them:

\begin{lstlisting}
data:
  data_paths:
    - data/MP
  potcar_dir: /opt/vasp/potpaw_PBE     # <dir>/Ag/POTCAR, <dir>/Pt/POTCAR, ...
  resolution: 32
\end{lstlisting}

\file{.gz} and \file{.Z} entries are read directly, and a flat
\file{POTCAR.Ag} / \file{Ag.POTCAR} layout is recognised too. A run that has
its own \file{POTCAR} ignores the library entirely --- it is a fallback, not an
override, and it also rescues local runs whose \file{POTCAR} was stripped for
licensing.

\begin{center}
\begin{tabular}{@{}lll@{}}
\toprule
Pseudopotentials & $\vext$ & residual vs.\ a reference \file{EXTCAR} \\
\midrule
available & tabulated local pseudopotential & $2\times10^{-5}$ relative $L^2$ \\
none & Gaussian pseudo-ion model & order $10^{-1}$ \\
\bottomrule
\end{tabular}
\end{center}

Measured on the Ag--Pt--Pt set, the two constructions differ \emph{from each
other} by $0.38$ relative $L^2$. They are different fields, not different
roundings of one; the Technical Guide shows why no analytic form factor can
close that gap.

\begin{pwarn}
\textbf{Use the library that generated the data.} The Materials Project uses
the VASP PBE set (\opt{PAW\_PBE}). A \opt{PAW\_LDA} library would build a
potential the densities were never computed from --- worse than the Gaussian
fallback, because it looks exact.
\end{pwarn}

Missing or unreadable entries are not fatal. Each such species warns once and
falls back to the Gaussian model for the structures containing it, so a library
covering four elements of five still buys the exact potential for everything
that does not involve the fifth. The run log names the construction per source.

\subsection{What the cache records, and \opt{strict\_potcar}}
\label{sec:strict-potcar}

The cache \emph{directory name} carries \opt{potcar} when a library was
configured, and until 26.9.3 that was the whole of the record --- which is to
say it recorded what was \emph{asked for} and nothing about what was
\emph{got}. The fingerprint had the same gap.

It cost a measurement. A control run on Santos Dumont landed on a node where
the library's filesystem was not mounted, warned six times on stderr, trained
against analytic pseudo-ion potentials, and wrote a
\file{cache\_fingerprint.json} identical to one built from real
pseudopotentials --- in a directory called \file{res32\_potcar}. Every later
run reused it silently, on mounted nodes too, and the warning never came back.

The fingerprint now carries \opt{potcar\_source}, per element:

\begin{lstlisting}
{ "potcar_dir": "/prj/.../POTCARs",
  "potcar_source": {"Ag": "gaussian", "Pt": "library"} }
\end{lstlisting}

so a cache built during an outage no longer matches a run that can read the
library --- it rebuilds rather than being reused --- and the resolved mode is
in the run log beside the rest of the cache line. It is \opt{null} when no
library is configured, so a purely Gaussian dataset's fingerprint is unchanged.

\opt{data.strict\_potcar: true} turns the fallback into an error naming the
elements the library did not serve. Off by default, and meant for queues: the
failure shape is \opt{training.strict\_device}'s --- a job that holds its
allocation and quietly computes something other than what was configured.

Without a library the map learned is \emph{model potential} $\to$ \emph{DFT
density}: well-posed and self-consistent, since inference builds the very same
potential, but not VASP's $\vext$. Say so wherever such a model's numbers are
quoted.

\section{Training on several sources at once}

\opt{data.data\_paths} is a list of directories, and \textbf{every entry has
the same shape}: subdirectories, one per material, each holding that material's
volumetric files. Nothing in the config says which is which, because on disk
they are the same kind of thing:

\begin{lstlisting}
data:
  data_paths:
    - data/vasp/structures  # structure_0000/, structure_0001/, ...
    - data/MP               # mp-124/, mp-81/, ...
    - data/cache/res32      # a cache from an earlier run
  potcar_dir: /opt/vasp/potpaw_PBE
\end{lstlisting}

What differs is the \emph{content} of a material's directory --- inputs beside
the density mean $\vext$ is built from them, a density alone means it is built
from the structure the \file{CHGCAR} carries in its own header --- and content
is read, not declared.

Everything found is pooled into one training set, and identifiers that collide
between directories are prefixed with their own so no material is silently
dropped. \file{TAUCAR} is optional \textbf{per material}: one run may have
written one while its neighbour did not, and the one that did still reaches
\opt{chg2tau} while both reach \opt{ext2chg}.

\begin{pwarn}
Without \opt{potcar\_dir}, mixing materials that carry their own
pseudopotentials with materials that carry only a density mixes \emph{two
definitions of} $\vext$ under one name, and the operator spends
capacity reconciling them. Poraqu\^e warns when it detects this. Setting
\opt{potcar\_dir} makes both sources use the tabulated construction, which
resolves the warning correctly rather than merely silencing it.
\end{pwarn}
